\begin{exercice*}[Extrait du DNB]
    Il y a dans une urne douze boules indiscernables au toucher, numérotées de 1 à 12. On veut tirer une boule au hasard.

    \begin{enumerate}
        \item Est-il plus probable d'obtenir un numéro pair ou bien un multiple de 3 ?
        \item Quelle est la probabilité d'obtenir un numéro inférieur à 20 ?
        \item On enlève de l'urne toutes les boules dont le numéro est un diviseur de 6. On veut à nouveau tirer une boule au hasard. Explique pourquoi la probabilité d'obtenir un numéro qui soit un nombre premier est alors 0,375.
    \end{enumerate}
\end{exercice*}
\begin{corrige}
    %\setcounter{partie}{0} % Pour s'assurer que le compteur de \partie est à zéro dans les corrigés
    %\phantom{rrr}    
    \begin{enumerate}
        \item Il y a 6 numéros pairs (2, 4, 6, 8, 10 et 12) et 4 numéros multiples de 3 (3, 6, 9 et 12). Les numéros pairs sont plus nombreux que les numéros multiples de 3, donc il est plus probable d'obtenir un numéro pair.
        \item Tous les numéros sont inférieurs à 20, donc la probabilité d'obtenir un numéro inférieur à 20 est de $\frac{12}{12} = 1$.
        \item Les diviseurs de 6 sont : 1, 2, 3 et 6. On enlève donc les boules numérotées de ces nombres. Il reste alors les boules numérotées : 4, 5, 7, 8, 9, 10, 11 et 12. Parmi ces boules restantes, il y a trois numéros premiers : 5, 7 et 11. La probabilité d'obtenir un numéro qui soit un nombre premier est alors de $\frac{3}{8} = 0,375$.
    \end{enumerate}
\end{corrige}

